

\def\tituloPortadaVal{\titulo}
\newcommand{\tituloPortada}[1]{%
\def\tituloPortadaVal{#1}
}

% Autor
\def\autorPortadaVal{\autor} \newcommand{\autorPortada}[1]{%
  \def\autorPortadaVal{#1} }

% Tutor
\def\directorPortadaVal{Tutor no definido. Usa
  \texttt{$\backslash$tutorPortada}} \newcommand{\directorPortada}[1]{%
  \def\directorPortadaVal{#1} }

% Fecha de publicacion
\def\fechaPublicacionVal{} \newcommand{\fechaPublicacion}[1]{%
  \def\fechaPublicacionVal{#1} }

% Imagen de la portada (escudo)
\def\imagenPortadaVal{Imagenes/Vectorial/Todo} \newcommand{\imagenPortada}[1]{%
  \def\imagenPortadaVal{#1} } \def\escalaImagenPortadaVal{1.0}
\newcommand{\escalaImagenPortada}[1]{%
  \def\escalaImagenPortadaVal{#1} }
% Tipo de documento (TESIS, MANUAL, ...)
\def\tipoDocumentoVal{TESIS DOCTORAL} \newcommand{\tipoDocumento}[1]{%
  \def\tipoDocumentoVal{#1} }

% Instituci_n
\def\institucionVal{} \newcommand{\institucion}[1]{%
  \def\institucionVal{#1} }

% Primer subt_tulo de la segunda portada
\def\textoPrimerSubtituloPortadaVal{%
  \textit{Memoria que presenta para optar al master en ingeniería informática}
  \leavevmode \\[0.3em]%
  \textbf{\autorPortadaVal}\\[1.2em]
   %Calificación: \bb{Matrícula de Honor}
} \newcommand{\textoPrimerSubtituloPortada}[1]{%
  \def\textoPrimerSubtituloPortadaVal{#1} }

% Segundo subt_tulo de la segunda portada
\def\textoSegundoSubtituloPortadaVal{%
  \textit{Directed by} \leavevmode \\[0.3em]
  \textbf{\directorPortadaVal} \\[0.5em]
   Calificación: 10 - \bb{Matrícula de Honor}
} \newcommand{\textoSegundoSubtituloPortada}[1]{%
  \def\textoSegundoSubtituloPortadaVal{#1} }

% ISBN
\newcommand{\isbn}[1]{%
  \def\isbnVal{#1} }

% Copyright
\newcommand{\copyrightInfo}[1]{%
  \def\copyrightInfoVal{#1} }

% Cr_ditos a TeXiS
\newcommand{\noTeXiSCredits}{%
  \def\noTeXiSCreditsVal{} }

% Explicaci_n sobre impresi_n a doble cara
\newcommand{\explicacionDobleCara}{%
  \def\explicacionDobleCaraVal{} }
%%%
% Configuraci_n terminada
%%%


%%%
%% COMANDO PARA CREAR LAS PORTADAS. CONTIENE TODO EL C_DIGO LaTeX
%%%
\newcommand{\makeCover}{%


  % Ponemos el marcador en el PDF
  \ifpdf \pdfbookmark{Página de Título}{titulo} \fi

%
  % M_RGENES
%
  % La maquetaci_n de las p_ginas en LaTeX es bastante complicada en lo que se
  % refiere a los m_rgenes. Por ejemplo, _siempre_ hay un desplazamiento de una
  % pulgada hacia la derecha y hacia abajo, porque por razones de maquetaci_n se
  % cree que es necesario ese espacio. Si quieres (en una hoja impar, que son de
  % las _nicas que me he preocupado) escribir por encima de la primera pulgada,
  % o hacia la izquierda, tienes que usar "m_rgenes" negativos.
%
  % Para poder seguir esto un poco, lo mejor es que ejecutes \layout con el
  % paquete layout cargado, para ver una imagen...
%
  % Los valores importantes en lo que se refiere al margen (horizontal, que es
  % del _nico que me he preocupado) son:
%
  % - \hoffset : desplazamiento horizontal del "eje" de coordenadas. A este
  % valor hay que sumarle, irremediablemente, una pulgada. El valor por defecto
  % es 0 pt - \oddsidemargin : "margen" izquierdo (en las p_ginas impares). El
  % texto principal de los p_rrafos comenzar_ en esa posici_n (es decir, 1
  % pulgada + \hoffset + \oddsidemargin). El valor por defecto es 22 pt -
  % \textwidth : longitud del texto (de los p_rrafos). El valor por defecto son
  % 360 pt. - \marginparsep : separador de la parte derecha del texto principal
  % y el espacio a anotaciones en el margen (notas marginales). El valor por
  % defecto es 7 pt. - \marginparwidth : ancho de la secci_n de notas
  % marginales. El valor por defecto es 106 pt.
%
  % Fijate que las notas marginales NO necesariamente llegar_n hasta el final
  % del folio. La separaci_n entre el extremo derecho de las notas al margen y
  % el final del folio en realidad depender_ del tama_o
  % de% _ste; no se especifica de ninguna manera. Esto significa que si
  % quieres ajustar exactamente el margen derecho tienes que echar cuentas con
  % respecto al tama_o del papel, que se mantiene en \paperwidth.
%
  % En realidad, no s_ de qu_ manera pero todos los valores est_n relacionados
  % de alg_n modo, y un cambio en uno afecta a los dem_s de maneras bastante
  % inveros_miles. Adem_s, _s_lo_ pueden cambiarse en el pre_mbulo (bueno... al
  % menos \textwidth, \oddsidemargin se puede cambiar en otros sitios, pero no
  % se pueden hacer muchas cosas s_lo con aquellos que se pueden cambiar...).
%
  % En general se recomienda que NO cambies los m_rgenes. Est_n elegidos por
  % especialistas que saben de maquetaci_n y que han estudiado en profundidad
  % las mejores organizaciones. Por ejemplo, los cambios que hagas pueden
  % alargar demasiado las l_neas, o dejarlas demasiado cortas. Pero bueno, si
  % aun as_ quieres cambiar los m_rgenes _a nivel global_ es preferible que uses
  % el paquete geometry, cuya inclusi_n recibe los cent_metros de m_rgen que
  % quieres en cada lado, y _l se encarga de hacer las cuentas para que queden
  % as_, porque tocar tanto valor es un infierno.
%
  % Si quieres cambiar los m_rgenes moment_neamente, entonces lo mejor es hacer
  % uso de un entorno tipo "lista" que permite tocar algunos contadores para
  % ajustar las posiciones de los p_rrafos. Eso es precisamente lo que hace el
  % entorno cambiamargen definido en TeXiS.sty
%
  % El problema es que esos cambios son _relativos_ a los m_rgenes oficiales. Si
  % quieres hacer un cambio dr_stico (como el que necesitamos en la portada,
  % para que quede centrada), entonces es necesario echar cuentas con los
  % contadores anteriores para realizar el desplazamiento adecuado en cada lado
  % para conseguirlo.
%
  % Para resumir, los valores usados en los m_rgenes eran: 1 pulgada + hoffset +
  % oddsidemargin + + textwidth + + marginparsep + marginparwidth + AJUSTE =
  % paperwidth
%
  % Lo que necesitamos en crear un entorno cambiamargen pasando los valores
  % adecuados para que quede centrado. Para eso, hay que hacer cuentas, y eso es
  % un tanto infierno en LaTeX a no ser que se incluya el paquete calc que
  % permite notaci_n infija de operadores. Por tanto, para que esto funcione hay
  % que incluirlo.
%
  % Para aclararnos, vamos a crear un par de longitudes para hacer las cuentas
  % poco a poco. Adem_s, lo primero es asumir que queremos _eliminar todos los
  % m_rgenes_ y que los p_rrafos lleguen totalmente de lado a lado. Para eso, en
  % la izquierda tendremos que restar la suma de los tres primeros valores (1
  % pulgada, \hoffset y \oddsidemargin).

  \newlength{\cambioIzquierdo} \setlength{\cambioIzquierdo}{1in + \hoffset +
    \oddsidemargin}
  % Si te falla aqu_, incluye el paquete calc en el pre_mbulo. 1in = 1 pulgada =
  % 72.27 pt

  % En la parte derecha hay que restar el "margen" visible total, que es el
  % tama_o de la hoja restando el espacio hasta la izquierda del p_rrafo y su
  % ancho. Aprovechamos que el primer espacio lo tenemos en \cambioIzquierdo.
  \newlength{\cambioDerecho} \setlength{\cambioDerecho}{\cambioIzquierdo +
    \textwidth} \setlength{\cambioDerecho}{\paperwidth - \cambioDerecho}

  % Ya casi est_. Si hicieramos
%
% \begin{cambiamargen}{-\cambioIzquierdo}{-\cambioDerecho}
%   ...
% \end{cambiamargen}
%
% tendriamos p_rrafos que van de extremo a extremo de la hoja. Como eso es una
% exageraci_n, restamos a cada longitud el m_rgen real que queremos dejar.
  \newlength{\margenPortada} \setlength{\margenPortada}{3.5cm}

  \setlength{\cambioIzquierdo}{\cambioIzquierdo - \margenPortada}
  \setlength{\cambioDerecho}{\cambioDerecho - \margenPortada}


%%%
  % Portada
%%%


  % P_gina sin cabeceras
  \thispagestyle{empty}

  % \begin{cambiamargen}{\cambioIzquierdo}{-\cambioDerecho}

  \newgeometry{inner=2cm,outer=3cm,top=2cm}

  % En la primera hoja no se entiende de p_ginas pares e impares
  \newlength{\evensidemarginOriginal}
  \setlength{\evensidemarginOriginal}{\evensidemargin}

  \newlength{\oddsidemarginOriginal}
  \setlength{\oddsidemarginOriginal}{\oddsidemargin}

  \setlength{\evensidemargin}{0cm} \setlength{\oddsidemargin}{0cm}

  % Comienza el tex...
  \begin{huge}
    \mbox{ } \vfill %

    % Cuidado: Ajustarlo si el t_tulo cambia...
    \newlength{\longTitulo} \settowidth{\longTitulo}{%
      \textbf{Arquitectura y metodología para el}}
    \begin{center}
      \rule{\longTitulo}{.5mm}\leavevmode \\\relax \vskip 1cm
      \textbf{\tituloPortadaVal} \vskip 0.7cm
      \rule{\longTitulo}{.5mm}\leavevmode \\\relax
    \end{center}
  \end{huge}

  \vfill %

\begin{center}
  \includegraphics[scale=\escalaImagenPortadaVal]{\imagenPortadaVal}
\end{center}

\vfill %

\begin{center}
  {\Large \textbf{\tipoDocumentoVal}}
\end{center}

\vfill

\begin{large}

\begin{center}
  \textbf{\autorPortadaVal}%\leavevmode \\[0.3cm]


  \leavevmode \\\relax \mbox{ }% \leavevmode \\% \mbox{ } \leavevmode \\
  \textoSegundoSubtituloPortadaVal \leavevmode \\[0.3]
  \leavevmode \\\relax \mbox{ }% \leavevmode \\ \mbox{ } \leavevmode \\

  \textbf{\institucionVal}\leavevmode \\[1em]
  \textbf{\fechaPublicacionVal}

\end{center}

\end{large}

\vfill %

% \end{cambiamargen}

\restoregeometry


\newpage

% P_gina que en la cara de detr_s del folio de la portada. Ponemos que el
% documento esta maquetado con TeXiS y la aclaraci_n de que debe imprimirse a
% doble cara.
\thispagestyle{empty} \mbox{ } \vfill%space*{4cm}
\begin{small}
  \begin{center}
    \ifx\noTeXiSCreditsVal\undefined Document made with \texis\ v.\texisVer.
    modified by Viktor Jacynycz García \else \mbox{ } \fi
  \end{center}
\end{small}
\vspace*{2cm}
\begin{small}
  \begin{center}
    \ifx\explicacionDobleCaraVal\undefined \mbox{ } \else
    \noindent This document is prepared for duble-side printing. \fi
  \end{center}
\end{small}

%%%
% Segunda portada
%%%

\newpage

\thispagestyle{empty}

\mbox{ }

\begin{Huge}
  \begin{center}
    \tituloPortadaVal
  \end{center}
\end{Huge}

\vfill %

\begin{large}
  \begin{center}
    \textoPrimerSubtituloPortadaVal\ \leavevmode \\\relax \mbox{ } \leavevmode
    \\\relax \mbox{ } \leavevmode \\\relax
    \textoSegundoSubtituloPortadaVal \leavevmode \\[0.3em]
  \end{center}
\end{large}

\vfill %

\begin{large}
  \begin{center}
    \textbf{\institucionVal}\leavevmode \\[0.2em]
    \mbox{ } \leavevmode \\\relax \textbf{\fechaPublicacionVal}
  \end{center}
\end{large}


\newpage
\thispagestyle{empty} \mbox{ }

% Informacion del ISBN y copyright
\vskip 6cm
\begin{center}
\textbf{Attribution-NonCommercial-ShareAlike 4.0 International}

  \emph{CC BY-NC-SA 4.0}
\end{center}
\vskip 5mm
You are free to:
\begin{itemize}


\itbf{Share} — copy and redistribute the material in any medium or
format \itbf{Adapt} — remix, transform, and build upon the material
\end{itemize}

The licensor
  cannot revoke these freedoms as long as you follow
  the following terms:

  \begin{itemize}
    \itbf{Attribution} — You must give appropriate credit, provide a
  link to the license, and indicate if changes were made. You may do so in any
  reasonable manner, but not in any way that suggests the licensor endorses you
  or your use.

  \itbf{NonCommercial} — You may not use the material for commercial purposes.

  \itbf{ShareAlike} — If you remix, transform, or build upon the material, you must
  distribute your contributions under the same license as the original.

  \itbf{No additional restrictions} — You may not apply legal terms or technological
  measures that legally restrict others from doing anything the license permits.
  \end{itemize}
\begin{center}

  \includegraphics{TeXiS/by-nc-sa.png}\\%
  https://creativecommons.org/licenses/by-nc-sa/4.0/\\%
  Viktor Jacynycz García
\end{center}
\newpage
\thispagestyle{empty} \mbox{ } \vskip 5cm El abajo firmante, matriculado en el
Máster en Ingeniería Informática de la Facultad de Informática, autoriza a la
Universidad Complutense de Madrid (UCM) a difundir y utilizar con fines
académicos, no comerciales y mencionando expresamente a su autor el presente
Trabajo Fin de Máster: \emph{\tituloPortadaVal}, realizado durante el curso
académico 2017-2018 bajo la dirección de Samer Hassan y Antonio Sanchez
Ruiz-Granados en el Departamento de ISIA, y a la Biblioteca de la UCM a
depositarlo en el Archivo Institucional E-Prints Complutense con el objeto de
incrementar la difusión, uso e impacto del trabajo en internet y garantizar su
preservación y acceso a largo plazo.

\mbox{ } \vskip 5cm
\begin{center}
  Fdo. Viktor Jacynycz García
\end{center}
}% \newcommand{\makeCover}

% Variable local para emacs, para que encuentre el fichero maestro de
% compilaci_n
%%%
%%% Local Variables:
%%% mode: latex
%%% TeX-master: "../Tesis.tex"
%%% End:
